% !PRINTED 191
\weberpage{191}
%<<CONT

\weq{14}{(\alpha'_h\beta'_k) = r(\alpha_h\beta_k),}

according to which the determinants $(\alpha_h\beta_k)$ are to be regarded as irrational \gesperrt{invariants}. We now wish to set out how rational invariants and covariants can be formed from them.

For this purpose it is best to consider the linear transformation in the inhomogeneous form

\weq{15}{x = \frac{\alpha x' + \beta}{\gamma x' + \delta}}

and to apply it to the function $f(x)$ in (1). It then follows, if

\[
F(x') = a'_0 x'^n + a'_1 x'^{n-1} + a'_2 x'^{n-2} + \cdots + a'_n
\]

is set, that

\weq{16}{F(x') = (\gamma x' + \delta)^n f(x).}

We also write this equation thus

% [sic] The print omits the otherwise required plus sign between the a_2/x^2 term and the ellipsis.
\[
F(x') = (\alpha x' + \beta)^n \left(a_0 + a_1 \frac{1}{x} + a_2 \frac{1}{x^2} \cdots\right),
\]

from which, if one sets $x' = -\delta : \gamma$, and hence $x = \infty$,

\[
a_0 r^n = (-\gamma)^n F\left(-\frac{\delta}{\gamma}\right),
\]

or, if $\alpha'_1, \alpha'_2, \ldots \alpha'_n$ depend on $\alpha_1, \alpha_2, \ldots \alpha_n$ in the same way as $x'$ does on $x$, so that

\[
F(x') = a'_0 (x' - \alpha'_1)(x' - \alpha'_2) \ldots (x' - \alpha'_n)
\]

is set,

\weq{17}{r^n a_0 = a'_0(\gamma\alpha'_1 + \delta)(\gamma\alpha'_2 + \delta) \ldots (\gamma\alpha'_n + \delta)}

follows.

Now by the transformation (15) we have

\weq{18}{\begin{aligned}
x - \alpha_i &= \frac{r(x' - \alpha'_i)}{(\gamma x' + \delta)(\gamma\alpha'_i + \delta)} \\
\alpha_i - \alpha_k &= \frac{r(\alpha'_i - \alpha'_k)}{(\gamma\alpha'_i + \delta)(\gamma\alpha'_k + \delta)}.
\end{aligned}}

We now form a product $P$ out of any of the factors

\[
\begin{aligned}
&(0, 1), (0, 2) \ldots (0, n) \\
&(1, 2), (1, 3) \ldots (n - 1, n),
\end{aligned}
\]

but in such a way that each of the digits $1, 2 \ldots n$ occurs in it equally often in all, say $\mu$ times. The index $0$, that is, the variable $x$, may occur $\nu$ times, and the total number of the
%CONT>>
